% Independent pencil argument by /root; checked by /root/cycles and /root/frobenius.
\paragraph{C.2e (Proved under the stated field-extension hypothesis): $A'_H$ implies $A_H$.}
Fix an algebraically closed field $k$ of arbitrary characteristic and a
prime $\ell\ne\operatorname{char}(k)$. We use $\ell$-adic cohomology,
with rational cycle coefficients, and write $A^r_H$ for the image of
$\mathrm{CH}^r(-)_{\mathbb Q}$ in $H^{2r}(-,\mathbb Q_\ell(r))$.
Assume $A'_H$ for all smooth projective varieties and smooth ample
hyperplane sections over every algebraically closed extension of $k$.
Then $A_H$ holds over all those extensions. In particular the universal
implication in the letter is valid. The extension hypothesis is used
at the geometric generic fibre of a pencil; we do not replace it by
an unsupported assertion about specializing algebraic preimages.

\begin{proof}
We prove $A_H(X,i)$ first by induction on $i$, and, for fixed $i$,
by induction on $n=\dim X\ge2i$. The induction includes all the
specified base fields simultaneously. For $i=0$, both cycle spaces are
one-dimensional by degree, and $L^n$ has positive nonzero degree.
For any $i$, the case $n=2i$ is the identity. Replacing an ample line
bundle by a sufficiently high very ample power does not change the
assertion: it multiplies each relevant power of $L$ by a nonzero
rational scalar. We therefore work with a projective embedding.

Let $n\ge2i+1$, set $\kappa=n-2i$, and take
$z\in A_H^{n-i}(X)$. Choose a pencil whose base locus $B$ is a smooth
codimension-two complete intersection and whose generic member is
smooth. Bertini supplies this over the algebraically closed base.
No assertion about ordinary double points of the other members is
needed. Write
\[
 \pi:\widetilde X=\operatorname{Bl}_B X\longrightarrow X,
 \qquad f:\widetilde X\longrightarrow\mathbb P^1,
 \qquad K=k(t),
\]
and let $Y_\eta$ be the generic fibre, viewed as a hyperplane section of
$X_K$. Over $\overline K$, the hypothesis $A'_H$ gives a class
$b\in A_H^{n-i-1}(Y_{\overline\eta})$ whose Gysin image is $z$;
its degree lies in the required range since
$2(n-i-1)\ge n-1$. The induction on dimension gives
\[
 b=L_Y^{\kappa-1}\gamma,
 \qquad \gamma\in A_H^i(Y_{\overline\eta}).
\]
Here $n-1\ge2i$, including the critical case $n=2i+1$, where this
induction step is just the identity map. Thus
\begin{equation}\label{eq:C2-generic-push}
             z=L^{\kappa-1}(i_\eta)_*\gamma.
\end{equation}

A rational cycle representing $\gamma$ is defined over a finite
extension $K'/K$. Transfer divided by $[K':K]$ gives a rational cycle
on $Y_\eta$ for which \eqref{eq:C2-generic-push} still holds: the
class $z$ is fixed under base-field automorphisms, and pushforward,
cycle class and transfer commute. Purely inseparable factors multiply
pushforward by their degrees and induce the usual isomorphism on
prime-to-characteristic cohomology, so cause no exception. We keep
$\gamma$ for this averaged cycle and close it in $\widetilde X$.
The Chow blow-up formula, with the explicit maps of
\citet[Theorem 2.8(3)]{cyJiang2023}, writes its closure as
\[
 \Gamma=\pi^*a+j_*q^*\beta,
 \qquad a\in\mathrm{CH}^i(X)_{\mathbb Q},\quad
 \beta\in\mathrm{CH}^{i-1}(B)_{\mathbb Q},
\]
where $j:E\hookrightarrow\widetilde X$ is the exceptional divisor
and $q:E\to B$ its projection. The pushforward of the restriction to
one pencil member is
\begin{equation}\label{eq:C2-exceptional}
                  (i_\eta)_*\gamma=L a+(i_B)_*\beta.
\end{equation}
For clarity, the fibre class is $F=\pi^*L-E$.
We have $\pi_*(F\pi^*a)=La$, since $\pi_*E=0$.
Moreover $F|_E$ has degree one on each exceptional $\mathbb P^1$,
so $q_*(F|_E)=1$ and
$\pi_*(Fj_*q^*\beta)=(i_B)_*\beta$.
This proves both the sign and the coefficient in
\eqref{eq:C2-exceptional}.

Weak Lefschetz for the two hyperplane sections identifies
\[
 H^{2i-2}(X,\mathbb Q_\ell(i-1))
       \xrightarrow{\sim}
 H^{2i-2}(B,\mathbb Q_\ell(i-1)),
\]
since $2i-2<n-2=\dim B$. Let $b_0$ be the preimage of the cycle
class of $\beta$. The projection formula gives
$(i_B)_*\beta=L^2b_0$. The class
\[
 L^{\kappa+2}b_0=L^\kappa(i_B)_*\beta
\]
is algebraic. The already proved induction step with codimension
$i-1$ on $X$ makes $L^{\kappa+2}$ surjective from
$A_H^{i-1}(X)$ onto $A_H^{n-i+1}(X)$. Hard Lefschetz on the full
cohomology makes it injective. Consequently $b_0$ itself is algebraic.
Combining the two displayed pushforward formulas now gives
\[
                 z=L^\kappa(a+L b_0),
\]
which proves surjectivity. Hard Lefschetz supplies injectivity and
completes both inductions.
\end{proof}

This proof explains the role of a nonclosed intermediate field in the
letter. The actual cycle to be closed lives over $k(t)$; existence of
its preimage is applied over $\overline{k(t)}$, then finite transfer
returns it to $k(t)$. If one assumes $A'_H$ only for varieties over a
single fixed algebraically closed field, extension stability is an
additional premise in this argument. An assertion for every closed
member of a pencil is not substituted for the needed generic one.

\paragraph{C.2f (Proved): the converse and the exact affine-vanishing gate.}
For a smooth ample hyperplane section $i:Y\hookrightarrow X$,
$A_H(X)$ already implies $A'_H(X,Y)$. Indeed, for $2r\ge n-1$,
write an algebraic class $z\in A_H^{r+1}(X)$ as
\[
 z=L^{2r+2-n}a,\qquad a\in A_H^{n-r-1}(X),
\]
using $A_H(X,n-r-1)$. Since $2r+1-n\ge0$, the algebraic class
$L_Y^{2r+1-n}i^*a$ lies in $A_H^r(Y)$, and its Gysin image is $z$
by the projection formula. Thus, with the extension-stable universal
quantifiers of C.2e, $A_H$ and $A'_H$ are equivalent.

Since $A''_H$ is already a theorem by C.2b, the universal implication
$A''_H\Rightarrow A'_H$ would prove this standard conjecture $A$;
conversely that standard conjecture proves the implication. Together
with the product-level equivalences in
\citet[Corollary 2.14]{cyKleiman1968}, its exact remaining gate is the
universal Lefschetz standard conjecture. This is a reduction, not a
refutation of the implication and not a proof of that conjecture.

\paragraph{C.2g (Proved): reduction of $A'_H$ to the critical dimension.}
It suffices in the universal $A'_H$ hypothesis to assume the cases
$n=2r+1$. The following proof also works with algebraic equivalence
in place of homological equivalence. We retain rational cycle
coefficients and the smooth very ample hyperplane setting of the
letter.

\begin{proof}
Let $Y\subset X$ have defining section $s_0$ of the very ample line
bundle $L$, let $n=\dim X$, and let $(n-1)/2\le r\le n-1$.
Set $d=2r+1-n$. If $d=0$, this is the assumed case. Otherwise put
\[
 T=X\times\mathbb P^d,\qquad
 H=\left\{\sum_{j=0}^d s_j(x)t_j=0\right\}\subset T,
\]
where $s_1,\ldots,s_d$ are general sections of $L$.
The line bundle defining $H$ is $L\boxtimes\mathcal O(1)$, which
is very ample. We choose the sections so that their common zero
locus with $s_0$ is smooth of codimension $d+1$ (or empty), by
successive Bertini on the already smooth $Y$. Then $H$ is smooth:
away from this base locus a derivative in a projective-coordinate
direction is nonzero; at a base point the independent differentials
$ds_0,\ldots,ds_d$ have no zero linear combination with coefficients
$[t_0:\cdots:t_d]$. Now $\dim T=2r+1$, so the critical case applies.

Given $z\in A_H^{r+1}(X)$, it supplies an actual rational cycle
$\gamma\in\mathrm{CH}^r(H)_{\mathbb Q}$ whose pushforward has
cohomology class $\operatorname{pr}_X^*z$. The Chow projective-bundle
formula gives unique classes
\[
 (i_H)_*\gamma=\sum_{j=0}^d a_j h^j,
 \qquad a_j\in\mathrm{CH}^{r+1-j}(X)_{\mathbb Q},
 \quad h=c_1(\mathcal O_{\mathbb P^d}(1)).
\]
Its cohomological counterpart implies $\operatorname{cl}(a_0)=z$.
We now extract a stronger fact about $a_0$ than cohomological
vanishing. On $U=X\setminus Y$, the section $s_0$ trivializes $L$,
and the triangular coordinate change
\[
 t_0\longmapsto t_0+\sum_{j=1}^d(s_j/s_0)t_j
\]
identifies $H_U$ with the coordinate hyperplane
$\mathbb P^{d-1}_U\subset\mathbb P^d_U$.
This regular projective-bundle automorphism preserves
$\mathcal O(1)$ and every base class. Its Gysin image on Chow groups
is therefore precisely the ideal generated by $h$.
Consequently $a_0|_U=0$ in the \emph{Chow group}, not only in
cohomology. Chow localization provides
$\beta\in\mathrm{CH}^r(Y)_{\mathbb Q}$ with $a_0=(i_Y)_*\beta$.
Taking classes proves $A'_H(X,Y)$.

With algebraic equivalence, the same projective-bundle decomposition
is respected by families of cycles, so its zeroth coefficient is
algebraically equivalent to $z$. The Chow-localization step is
unchanged. This proves the claimed critical-case reduction there
as well, without asserting that the homological pencil proof
C.2e automatically proves its algebraic-equivalence analogue.
\end{proof}

Thus the precise phrase in 14:45, ``il suffit sauf erreur de faire
$n=2i+1$'', has a proof in the stated hyperplane setting. Combined
with C.2e--f, the critical homological $A'_H$ cases already have the
full strength of the extension-stable universal standard conjecture
$A_H$. The product stabilization does not assume that an arbitrary
cohomologically vanishing cycle is rationally equivalent to zero.

\paragraph{C.2h (Proved): middle positivity and the positive trace form.}
For homological cycle classes, the universal middle-dimensional
positivity assertion $B_H$ in 14:45 implies the Hodge standard
positivity assertion in every algebraic degree. We give the reduction,
then specify exactly the additional Lefschetz hypothesis needed for
its consequences on correspondences.

Let $x\in A_H^r(X)$ be primitive, with $\dim X=n\ge2r$, and let
$i:Y\hookrightarrow X$ be a smooth complete intersection of dimension
$2r$ formed from the chosen hyperplanes. Restriction in degree $2r$
is injective by weak Lefschetz. Moreover $i^*x$ is primitive: the
Gysin map out of $H^{2r+2}(Y)$ is injective, by duality with weak
Lefschetz in degree $2r-2$, and
\[
 i_*(L_Y i^*x)=L^{n-2r+1}x=0.
\]
Thus middle positivity on $Y$ gives, for $x\ne0$,
\[
          (-1)^r\deg_X(L^{n-2r}x^2)
          =(-1)^r\deg_Y((i^*x)^2)>0.
\]
This proves the reduction in
\citet[Proposition 3.7]{cyKleiman1968}; its converse is its middle case.
No algebraicity of a Lefschetz inverse was used in this paragraph.

Now assume the Lefschetz standard conjecture $B_{\rm std}$ for the
varieties involved and Hodge standard positivity on their products.
The Lefschetz primitive projectors, inverse operators, and star
operators are algebraic. If $u$ is an algebraic cohomological
correspondence, its adjoint
\[
                  u^\dagger=*_{X}\,{}^t u\,*_{Y}
\]
(with the standard degree-wise Lefschetz signs) is therefore algebraic.
The ordinary trace $\operatorname{Tr}(u^\dagger u)$ is rational and
strictly positive for $u\ne0$. Here is the checked computation behind
this input, as in \citet[Theorem 3.11]{cyKleiman1968}.
Decompose $u$ by the algebraic orthogonal projectors onto Lefschetz
strings. The trace norm is the sum of the trace norms of these blocks.
Transporting each block to its primitive factors reduces to a
primitive class in $P^a(X)\otimes P^b(Y)$, where $a+b$ is even.
The product Lefschetz operator is $L_X\otimes1+1\otimes L_Y$;
expanding its power contributes the positive binomial factor
\[
                   \binom{n-a+m-b}{n-a},
                   \qquad n=\dim X,\quad m=\dim Y.
\]
The trace formula contributes $(-1)^a$, and exchanging the two
primitive factors contributes $(-1)^{ab}$. Their product is $+1$
because $a+b$ is even. Up to that positive binomial factor, the block
norm is therefore its primitive Hodge intersection norm on $X\times Y$,
which is positive by the stated hypothesis. Rationality follows from
the correspondence trace formula, an intersection number of rational
cycles. This verifies the hypotheses of the finite algebra arguments
below; positivity on just $X$, without the relevant products, is not
substituted for them.

The resulting $\mathbb Q$-algebra $E$ of algebraic endomorphisms is
finite dimensional, even before one invokes semisimplicity. To see
this directly, choose a maximal linearly independent list
$u_1,\ldots,u_s$ over the cohomology coefficient field in the image
algebra. The map
\[
 E\longrightarrow\mathbb Q^s,
 \qquad u\longmapsto
       (\operatorname{Tr}(u^\dagger u_1),\ldots,
                         \operatorname{Tr}(u^\dagger u_s))
\]
is injective. If all coordinates vanish, express $u$ in the coefficient-
field span of the $u_j$ to obtain
$\operatorname{Tr}(u^\dagger u)=0$, whence $u=0$.
The $u_j$ are also rationally independent, so the injection proves
$\dim_{\mathbb Q}E=s$. They are consequently a rational basis as well
as a basis after extension to the cohomology coefficient field.
In particular that scalar extension of the correspondence algebra
remains faithful. This proves finite dimensionality and a positive
definite rational trace form. Its extension to $E\otimes\mathbb R$ is positive definite:
a negative real direction would give a negative rational direction,
and a degenerate rational matrix would have a nonzero rational kernel.

\paragraph{C.2i (Proved under the preceding standard-conjecture hypotheses): semisimplicity and algebraic preimages.}
Under C.2h's hypotheses, $E$ is semisimple. Indeed its Jacobson radical
is preserved by the anti-involution $\dagger$. If $u$ belonged to that
radical, then $u^\dagger u$ would be nilpotent and would have zero trace
in cohomology. Positivity implies $u=0$. Thus the radical is zero.
This proves the finite-algebra assertion of
\citet[Corollary 3.12 and Lemma 3.13]{cyKleiman1968} in the present
trace realization. With the hypotheses on all products, the same
argument, or the numerical identification and
\citet[Theorem 1]{Jannsen1992}, gives the semisimplicity of the
idempotent-complete pure homological motive category.

We also prove the full image criterion in 14:46, without assuming
that a cohomological image automatically has an algebraic preimage.
Let $u:H^a(X)\to H^b(Y)$ be induced by a rational algebraic
correspondence, with the Tate twist matching its codimension.
Work in the algebra of blocks on the direct sum of these two
realizations. Put $x=u^\dagger u$. Left multiplication by $x$ on
this finite-dimensional real algebra is self-adjoint for the trace
inner product, so $x$ is semisimple. Its minimal polynomial has
rational coefficients and no repeated root. By the Chinese remainder
theorem there is $p(T)\in\mathbb Q[T]$ which is zero at $T=0$ and
agrees with $1/T$ at every nonzero root of that minimal polynomial.
Let
\[
                    e_0=1-p(x)x,\qquad v=p(x)u^\dagger.
\]
The element $e_0$ is a self-adjoint idempotent with $xe_0=0$.
Therefore
\[
 (ue_0)^\dagger(ue_0)=e_0xe_0=0,
\]
and trace positivity forces $ue_0=0$. It follows that
\[
                         uvu=u.
\]
Every operator used is algebraic. If an algebraic class $z$ lies
in the cohomological image of $u$, then $u(vz)=z$, and $vz$ is an
algebraic class. The converse is immediate. This proves exactly the
surjectivity-on-algebraic-images consequence, including maps that are
neither injective nor surjective. It also constructs the projectors
onto the image and kernel; the faithful realization verifies their
idempotent identities. The argument is the explicit pseudo-inverse
form of \citet[Corollary 3.14]{cyKleiman1968}.

These assertions require both algebraic Lefschetz operators and the
product positivity used above. Since universal middle $B_H$ gives
$H_{\rm std}$, the remaining implication ``$B_H$ alone implies $D$''
has, under weak and hard Lefschetz, the exact further requirement
$A_H$: C.2c gives $A_H+H_{\rm std}\Rightarrow D$ and
$D\Rightarrow A_H$. Over $\mathbb C$, where Hodge--Riemann positivity
already proves $B_H$, that requested implication has the full force
of the universal standard conjecture $A_H$. We have not proved it by
silently strengthening the meaning of the letter's $B$.

\paragraph{C.2j (Proved under explicit hypotheses): the conditional Weil consequences.}
Let $X/\mathbb F_q$ be smooth projective and geometrically connected
of dimension $n$. Use geometric Frobenius on $\ell$-adic cohomology,
$\ell\nmid q$, equivalently the pullback operator of the $q$-power
Frobenius endomorphism. Suppose $B_{\rm std}(X)$ and Hodge standard
positivity on $X\times X$ for the product polarization. Then its
operators $F_i$ on $H^i(X_{\overline{\mathbb F}_q},\mathbb Q_\ell)$
are semisimple, their characteristic polynomials have integer
coefficients, and every complex conjugate of every eigenvalue on
$H^i$ has absolute value $q^{i/2}$. We give the argument in order to
separate these conditional consequences from the already proved
Weil conjectures themselves.

First, $B_{\rm std}$ makes the cohomological projectors $\pi_i$
algebraic. Choose an integer $N>0$ clearing the denominators of a
cycle representing $\pi_i$. For every $m\ge1$, the trace formula gives
\[
             \operatorname{Tr}(F_i^m)\in N^{-1}\mathbb Z.
\]
Indeed the supertrace of the composition of $\pi_i$ with the integral
Frobenius correspondence of power $m$ is an intersection number;
only degree $i$ survives, so the ordinary trace differs by $(-1)^i$.
The same fixed $N$ works for every $m$. Newton's identities imply
$P_i(t)=\det(1-tF_i)\in\mathbb Q[t]$.

We include the bounded-denominator argument that removes every
prime denominator, including the characteristic prime. If a finite
multiset of algebraic numbers $\lambda_j$ has power sums in
$N^{-1}\mathbb Z$ for all positive powers, each $\lambda_j$ is an
algebraic integer. Otherwise, at some nonarchimedean place one of
them has absolute value greater than one. The power series
\[
 \sum_{m\ge1}\left(\sum_j\lambda_j^m\right)t^m
       =-\frac{tP'(t)}{P(t)},
 \qquad P(t)=\prod_j(1-\lambda_jt),
\]
has bounded coefficients at that place and is analytic on $|t|<1$.
The rational function on the right has a pole at the reciprocal
of such a $\lambda_j$, which lies in that disk. There is no cancellation:
a nonzero root of $P$ of multiplicity $e$ contributes a logarithmic
pole with nonzero residue proportional to $e$ in characteristic zero.
This contradicts the power-series identity. Thus all eigenvalues are
integral, and the rational elementary symmetric functions are integers.
This is the mechanism in \citet[Proposition 2.6]{cyKleiman1968}; it
requires no unproved removal of $p$-denominators.

Let $L$ be an ample class defined over $\mathbb F_q$, and extend the
scalar algebra from $\mathbb Q$ to $\mathbb Q(\sqrt q)\subset\mathbb R$.
The normalized degree-wise operator
\[
                   g=\sum_i q^{-i/2}F_i\pi_i
\]
is multiplicative, fixes $L$ because $F^*L=qL$, and acts as the
identity on top cohomology because the Frobenius morphism has degree
$q^n$. Poincar\'e duality therefore gives its transpose as $g^{-1}$.
Since $g$ commutes with $L$, with the primitive projectors, and with
Lefschetz star, its adjoint in C.2h is also $g^{-1}$:
\[
                             g^\dagger g=1.
\]
Left multiplication by $g$ on the finite real correspondence algebra
is an isometry for the positive trace inner product, since
\[
 \operatorname{Tr}((ga)^\dagger gb)
                   =\operatorname{Tr}(a^\dagger b).
\]
It is consequently diagonalizable over $\mathbb C$, with eigenvalues
on the unit circle. The left regular representation and the faithful
cohomological representation have the same minimal polynomial for an
algebra element: a polynomial in that element vanishes in one exactly
when it vanishes in the algebra. Hence the same semisimplicity and
absolute-value statement holds for $g$ in cohomology. On the invariant
left ideal $\pi_iE_{\mathbb R}$, left multiplication by $g$ equals
$q^{-i/2}$ times left multiplication by the rational algebra element
$F_i\pi_i$. The latter has a rational minimal polynomial; every complex
root occurs in its faithful left regular representation on that ideal.
The isometry therefore controls every complex conjugate of $F_i$, not
just one chosen embedding. Multiplying back by the positive real
number $q^{i/2}$ proves the claimed weights.
This supplies the linear algebra of
\citet[Lemma 4.6 and Theorem 4.7]{cyKleiman1968}, with the product
positivity and algebraic-projector premises visible.

Finally the fixed-point formula and the formal identity
\[
-\log\det(1-tF_i)=\sum_{m\ge1}\operatorname{Tr}(F_i^m)t^m/m
\]
give
\[
 Z(X,t)=\exp\left(\sum_{m\ge1}\#X(\mathbb F_{q^m})\frac{t^m}{m}\right)
       =\prod_{i=0}^{2n}P_i(t)^{(-1)^{i+1}}.
\]
Thus the numerator and denominator factors lie in $\mathbb Z[t]$;
$P_0(t)=1-t$, $P_{2n}(t)=1-q^nt$, and the zeros of $P_i$ have
absolute value $q^{-i/2}$. Poincar\'e duality pairs eigenvalues in
complementary degrees by $\alpha\leftrightarrow q^n/\alpha$, yielding
\[
 Z\left(X,\frac1{q^nt}\right)
        =\epsilon\,q^{n\chi/2}t^\chi Z(X,t),
 \qquad \chi=\sum_i(-1)^i\dim H^i(X),\quad \epsilon\in\{1,-1\}.
\]
More precisely, with
$\Delta=\prod_i\det(F_i)^{(-1)^{i+1}}$, the sign is
$\epsilon=(-1)^\chi q^{n\chi/2}\Delta$;
duality gives $\Delta^2=q^{-n\chi}$.
We retain this determinant sign instead of assuming that it is
fixed independently of Frobenius. Also $n\chi$ is even: when $n$
is odd, middle Poincar\'e duality is alternating and its dimension
is even, while the other terms occur in complementary pairs.
Rationality and the functional equation follow formally from the
trace formula and duality; the extra standard-conjecture input above
supplies the positive-involution proof of the absolute values and
semisimplicity. It is not a proof of the general standard conjectures.
